\section[Tổng kết]{Tổng kết}

\subsection{6 Kết luận}

\begin{frame}{Kết quả đạt được: kho chạy từ đầu đến cuối}
\muc{6 · Kết luận}
\HangSo[4]{\So{99.441}{đơn hàng}\hfill\So{16}{bảng chiều và fact}\hfill\So{17}{view mart}\hfill\So{136}{kiểm thử đạt}}
\begin{itemize}
  \item 10 chiều, 6 fact, 17 view mart; GMV 13,49 triệu BRL
\end{itemize}
\note{[S111] Tổng kết các con số chính của dự án: 10 chiều, 6 fact, 17 view mart; 99.441 đơn; GMV 13,49 triệu BRL; 136 kiểm thử đạt; bốn dashboard. Sơ đồ chòm sao (constellation) của kho nằm ở slide mở đầu. Tiêu đề trước đây: Kết quả: kho chạy được từ đầu đến cuối, 136 kiểm thử đạt, bốn dashboard}
\end{frame}

\begin{frame}{Đối chiếu với yêu cầu phân tích}
\muc{6 · Kết luận}
\only<1>{%
\begin{BangTC}{XlX}
\textbf{Yêu cầu (chủ điểm)} & \textbf{Dashboard} & \textbf{KPI đáp ứng} \\ \midrule
Bán hàng & DB1 & GMV, Orders, AOV, MoM, YoY, Pareto 80/20 \\
Xử lý đơn & DB2, DB3 & On-time rate, Avg lead days, Promised vs actual \\
Thanh toán & DB2 & cơ cấu thanh toán, số kỳ trả góp \\
\end{BangTC}
\KetLuan{Sáu chủ điểm ở mục 2.2 đều có dashboard và KPI tương ứng}\canxacnhan}
\only<2>{%
\begin{BangTC}{XlX}
\textbf{Yêu cầu (chủ điểm)} & \textbf{Dashboard} & \textbf{KPI đáp ứng} \\ \midrule
Trải nghiệm khách & DB2 & Avg review, \% negative, Repeat customer rate \\
Seller & DB4 & scorecard seller, GMV theo kênh \\
Thu hút seller & DB4 & MQLs, Won deals, MQL to won rate, Avg days to close \\
\end{BangTC}
\KetLuan{Sáu chủ điểm ở mục 2.2 đều có dashboard và KPI tương ứng}\canxacnhan}
\note{[P12] Đối chiếu từng yêu cầu ở mục 2.2 với dashboard và KPI đáp ứng. Cột đạt hay chưa nhóm xác nhận khi dashboard dựng xong (nhãn xác nhận là chữ giữ chỗ, cờ --final sẽ báo). Giám sát nạp dữ liệu (A11) nay là bảng ở mục 3.3, không còn dashboard.}
\end{frame}

\begin{frame}{Kiến thức thu được}
\muc{6 · Kết luận}
\begin{itemize}
  \item Thiết kế kho hình sao: grain, fact ba loại, SCD2, bus matrix
  \item ETL có kiểm thử: làm sạch ở ODS, đối soát số dòng và tổng tiền, chạy lại không đổi kết quả
  \item BI: mô hình ngữ nghĩa, DAX, drill-through, định dạng có điều kiện
  \item Mở rộng: lakehouse theo sự kiện và dữ liệu API công khai
\end{itemize}
\note{[P12] Kiến thức rút từ các phần 3, 4 và 5 của bài.}
\end{frame}

\begin{frame}{Bài học kỹ thuật: vấn đề, con số, cách xử lý}
\muc{6 · Kết luận}
\begin{BangTC}{XXX}
\textbf{Vấn đề} & \textbf{Con số} & \textbf{Cách xử lý} \\ \midrule
Một KPI có hai con số & GMV 13,49 so với 13,59 triệu; 380 so với 379 seller & chọn một định nghĩa, ghi cạnh con số \\
Điểm review đúng hạn sai & 4,21 rồi sửa thành 4,29 & sửa ở ETL (cột \texttt{is\_delivered}), không sửa ở BI \\
Khoá thay thế đổi giữa các lần dựng & hai bản dựng lệch khoá & chèn theo \texttt{ORDER BY} khoá tự nhiên \\
Dòng bị loại & 42 điểm toạ độ ngoài Brazil & ghi \texttt{etl\_log.rejected\_rows}, không xóa âm thầm \\
\end{BangTC}
\KetLuan{Đối soát với nguồn bắt được lỗi mà log không thấy}
\note{[P12] Bài học theo mạch vấn đề, con số, cách xử lý. Nguồn: các slide Giới hạn dữ liệu (2.3), ETL (3.3) và lakehouse (5.1). Mô hình chiều đáng dùng vì nhất quán; tốc độ đến từ materialized view và index.}
\end{frame}

\begin{frame}{Bài học làm việc nhóm}
\muc{6 · Kết luận}
\khungchen[0.3\textheight]{CHÈN: ba bài học làm việc nhóm, phân công, Git, khó khăn thực tế (Người 4 gom)}
\note{[P12] Khung chờ: Người 4 gom bài học làm việc nhóm từ cả nhóm.}
\end{frame}

\begin{frame}{Hạn chế và hướng phát triển}
\muc{6 · Kết luận}
\begin{itemize}
  \item Hạn chế: mô phỏng cập nhật chỉ thêm bản ghi mới; cơ chế thử lại mồ côi mới được chứng minh bằng test
  \item Hướng phát triển: lưu trạng thái trung gian của đơn
  \item Hướng phát triển: PostGIS cho tính toán địa lý
  \item Hướng phát triển: lịch chạy tự động
\end{itemize}
\note{[S113] Nhóm nói rõ giới hạn để kết quả đáng tin. Tiêu đề trước đây: Hạn chế còn lại và ba hướng phát triển: trạng thái đơn, PostGIS, lịch chạy}
\end{frame}

\begin{hustdarkframe}
\centering\vspace*{0.14\paperheight}
{\Huge\bfseries\color{white} CẢM ƠN QUÝ THẦY CÔ\\VÀ CÁC BẠN ĐÃ LẮNG NGHE!\par}\vspace{1em}
{\large\color{white!85} Hỏi đáp (Q\&A)\par}\vspace{0.8em}{\normalsize\color{white!85} Nhóm \TenNhom\ -- lớp \MaLop\ -- GVHD: \GiangVien\par}
\note{[S114] Nhóm sẵn sàng trả lời; các slide dự phòng B01–B10 dùng khi cần.}
\end{hustdarkframe}
